\documentclass[11pt,a4paper]{article}

% --- CJK 日本語対応 ---
\usepackage{CJKutf8}

% --- 数式・定理 ---
\usepackage{amsmath,amssymb,amsthm}

% --- 図表 ---
\usepackage{graphicx}
\usepackage{booktabs}
\usepackage{multirow}

% --- アルゴリズム ---
\usepackage[ruled,vlined]{algorithm2e}

% --- レイアウト ---
\usepackage[margin=25mm]{geometry}
\usepackage{enumitem}

% --- ハイパーリンク ---
\usepackage[hidelinks]{hyperref}

% --- 定理環境 ---
\newtheorem{theorem}{定理}
\newtheorem{lemma}[theorem]{補題}
\newtheorem{proposition}[theorem]{命題}
\newtheorem{corollary}[theorem]{系}
\newtheorem{definition}{定義}
\newtheorem{example}{例}
\newtheorem{remark}{注意}

\begin{document}
\begin{CJK}{UTF8}{min}

% ===================================================================
\title{最小記述長原理による密集区間要約}
\author{
  著者名 \\
  所属機関 \\
  \texttt{email@example.com}
}
\date{}
\maketitle

% ===================================================================
\begin{abstract}
頻出パターンマイニングは膨大な候補パターンを出力するため、
ユーザにとって結果の解釈が困難である。
本研究では、最小記述長（MDL: Minimum Description Length）原理を
密集区間構造に適用した新しいパターン要約手法
\textbf{Temporal Code Table（TCT）}を提案する。
TCTは各パターンに時間的有効区間を付与し、
\textbf{時間スコープ付き符号長}（Time-Scoped Code Length）により
区間内での出現を効率的に符号化する。
合成データ実験において、
提案手法は候補パターンの87--93\%を削減しつつ、
ベースライン対比70--88\%の圧縮率を達成した。
これは密集区間の時間的局所性がデータ圧縮に有効であることを示す。
\end{abstract}

\textbf{キーワード:}
最小記述長, パターン要約, 密集区間, KRIMP, コードテーブル


% ===================================================================
\section{はじめに}\label{sec:intro}

頻出アイテムセットマイニング~\cite{agrawal1994fast}は
データマイニングの基盤技術であるが、
低い最小サポート閾値の下では爆発的な数のパターンが出力される。
この問題に対し、Vreekenら~\cite{vreeken2011krimp}は
MDL原理に基づくKRIMPアルゴリズムを提案し、
データを最も簡潔に記述するパターンセットを選択する枠組みを確立した。

しかし、KRIMPおよびその後継手法（SLIM~\cite{smets2012slim}、
StreamKrimp~\cite{van2014streamkrimp}）は
パターンの時間的構造を考慮しない。
実世界のトランザクションデータでは、
特定のパターンが特定の時間区間でのみ頻出する
\textbf{密集区間}（dense interval）現象が広く観察される。
例えば、小売データにおける季節商品の共起や、
キャンペーン期間中の購買パターンの変化がこれに該当する。

本研究では、密集区間の時間的局所性を活用した
MDLベースのパターン要約手法を提案する。
具体的な貢献は以下の3点である:

\begin{enumerate}[label=(\arabic*)]
  \item \textbf{Temporal Code Table（TCT）}:
    各パターンに時間的有効区間を付与したコードテーブルを定義し、
    区間内での符号化効率を向上させる。
  \item \textbf{Time-Scoped Code Length}:
    パターンの区間内出現密度に基づく符号長を定義し、
    時間的に局所化されたパターンの記述コストを低減する。
  \item \textbf{Dense Interval Compression}:
    密集区間検出とMDL選択を統合した貪欲アルゴリズムにより、
    候補パターンの87--93\%を削減しつつ圧縮を達成する。
\end{enumerate}


% ===================================================================
\section{関連研究}\label{sec:related}

\subsection{MDLに基づくパターン選択}

MDL原理~\cite{rissanen1978modeling,grunwald2007minimum}は
モデルとデータの合計記述長を最小化する原理であり、
パターンマイニングにおけるモデル選択に広く応用されている。

\textbf{KRIMP}~\cite{vreeken2011krimp}は、
頻出パターンのコードテーブルを構築し、
Standard Cover Orderによる貪欲選択で
$L(\text{CT}) + L(D|\text{CT})$を最小化する。
\textbf{SLIM}~\cite{smets2012slim}はKRIMPを拡張し、
候補生成なしに直接コードテーブルを最適化する。
\textbf{StreamKrimp}~\cite{van2014streamkrimp}は
ストリーミング環境でのコードテーブル更新を扱うが、
パターンの時間的有効区間は明示的にモデル化しない。

\subsection{時間的パターンマイニング}

密集区間検出は、アイテムセットの出現が
特定の時間区間に集中する現象を捉える手法である。
スライディングウィンドウに基づく密集区間検出~\cite{apriori_window}は
効率的なアルゴリズムを提供するが、
出力パターン数の制御はユーザの閾値設定に依存する。

本研究は、密集区間検出の出力をMDL原理で要約するという
新しいアプローチを提案する。


% ===================================================================
\section{準備}\label{sec:prelim}

\subsection{記法}

$\mathcal{D} = \{T_1, T_2, \ldots, T_N\}$を
$N$個のトランザクションからなるデータベースとする。
各$T_t$はアイテム集合$\mathcal{I}$の部分集合である。
アイテムセット$X \subseteq \mathcal{I}$の
トランザクション$T_t$における出現を$X \subseteq T_t$と書く。

\begin{definition}[密集区間]
ウィンドウサイズ$W$と閾値$\theta$に対し、
アイテムセット$X$の\textbf{密集区間}$[s, e]$とは、
任意の$l \in [s, e]$について
$|\{t \in [l, l+W] : X \subseteq T_t\}| \geq \theta$
を満たす極大な連続区間である。
\end{definition}

\subsection{KRIMPの概要}

KRIMPでは、コードテーブル$\text{CT}$は
パターンと符号の対応表であり、
データの記述長は以下で定義される:
\begin{equation}\label{eq:krimp}
  L(\text{CT}, \mathcal{D}) = L(\text{CT}) + L(\mathcal{D} | \text{CT})
\end{equation}
ここで$L(\text{CT})$はコードテーブル自体の記述長、
$L(\mathcal{D}|\text{CT})$はコードテーブルを用いた
データの記述長である。

各パターン$X$の符号長は使用頻度に基づく:
\begin{equation}\label{eq:code_length}
  L(X | \text{CT}) = -\log_2 \frac{\text{usage}(X)}{\sum_{Y \in \text{CT}} \text{usage}(Y)}
\end{equation}


% ===================================================================
\section{提案手法}\label{sec:method}

\subsection{Temporal Code Table}

\begin{definition}[Temporal Code Table]
\textbf{Temporal Code Table（TCT）}は三つ組の集合
$\text{TCT} = \{(X_i, \mathcal{I}_i, c_i)\}_{i=1}^{k}$である。
ここで$X_i \subseteq \mathcal{I}$はパターン、
$\mathcal{I}_i = \{[s_{i,1}, e_{i,1}], \ldots, [s_{i,m_i}, e_{i,m_i}]\}$は
$X_i$の密集区間の集合、
$c_i$は使用回数である。
\end{definition}

TCTはKRIMPのコードテーブルを拡張し、
各パターンに時間的スコープ（有効区間）を付与する。
パターン$X_i$がトランザクション$T_t$のカバーに使用されるのは、
$t$がいずれかの密集区間$[s_{i,j}, e_{i,j}]$内にある場合に限る。

\subsection{Time-Scoped Code Length}

\begin{definition}[Time-Scoped Code Length]
パターン$X$の密集区間集合$\mathcal{I}_X$に対し、
\textbf{時間スコープ付き符号長}を以下で定義する:
\begin{equation}\label{eq:tscl}
  L_{\text{ts}}(X) =
    -\text{usage}_{\text{in}}(X) \cdot \log_2 \frac{\text{usage}_{\text{in}}(X)}{|\text{scope}(X)|}
    + \log_2 (N - |\text{scope}(X)|)
\end{equation}
ここで$\text{scope}(X) = \bigcup_{[s,e] \in \mathcal{I}_X} \{s, \ldots, e\}$は
有効範囲のトランザクション集合であり、
$\text{usage}_{\text{in}}(X)$は区間内での使用回数である。
\end{definition}

第1項は区間内での符号化コスト（密度が高いほど短い）、
第2項は区間外での出現に対するペナルティである。
密集パターンは区間内で高密度に出現するため、
通常のKRIMPよりも短い符号長を割り当てられる。

\subsection{Dense Interval Compression アルゴリズム}

提案アルゴリズムは以下の3段階で構成される:

\begin{enumerate}
  \item \textbf{候補生成}: Aprioriベースの密集区間検出により、
    密集区間を持つアイテムセットを列挙する。
  \item \textbf{Temporal Cover}: Standard Cover Orderに従い、
    各トランザクションを時間スコープ付きパターンでカバーする。
  \item \textbf{MDL選択}: 貪欲法でTCTにパターンを追加し、
    $L(\text{TCT}) + L(\mathcal{D}|\text{TCT})$が改善する場合にのみ採用する。
\end{enumerate}

\begin{algorithm}[t]
\SetAlgoLined
\KwIn{データベース$\mathcal{D}$, ウィンドウサイズ$W$, 閾値$\theta$}
\KwOut{Temporal Code Table $\text{TCT}^*$}

$\text{Candidates} \leftarrow \text{MineDensePatterns}(\mathcal{D}, W, \theta)$\;
$\text{TCT}^* \leftarrow \text{SCT}(\mathcal{D})$ \tcp*{Standard Code Table で初期化}
$L^* \leftarrow L(\text{TCT}^*) + L(\mathcal{D}|\text{TCT}^*)$\;

\ForEach{$(X, \mathcal{I}_X) \in \text{Candidates}$ in Standard Cover Order}{
  $\text{TCT}' \leftarrow \text{TCT}^* \cup \{(X, \mathcal{I}_X)\}$\;
  $L' \leftarrow L(\text{TCT}') + L(\mathcal{D}|\text{TCT}')$\;
  \If{$L' < L^*$}{
    $\text{TCT}^* \leftarrow \text{TCT}'$\;
    $L^* \leftarrow L'$\;
  }
}
\Return{$\text{TCT}^*$}\;
\caption{Dense Interval Compression}
\label{alg:dic}
\end{algorithm}

\begin{theorem}[圧縮保証]\label{thm:compression}
TCTはSCT以上の圧縮を保証する。
すなわち、
$L(\text{TCT}^*) + L(\mathcal{D}|\text{TCT}^*) \leq
 L(\text{SCT}) + L(\mathcal{D}|\text{SCT})$
が成り立つ。
\end{theorem}

\begin{proof}
アルゴリズム\ref{alg:dic}はSCTを初期解として開始し、
各候補追加ステップで記述長が改善する場合にのみ
パターンを採用する。
したがって、出力TCT$^*$の記述長は
初期解SCTの記述長以下である。
\end{proof}

\begin{theorem}[時間的局所性の利得]\label{thm:locality}
パターン$X$が区間$[s,e]$内でのみ出現し、
区間内密度$\rho = \text{usage}(X)/(e-s+1)$を持つとする。
TCTにおける$X$の符号長は、
KRIMPにおける$X$の符号長に対し、
\begin{equation}
  L_{\text{TCT}}(X) - L_{\text{KRIMP}}(X)
  \leq \text{usage}(X) \cdot \log_2 \frac{e-s+1}{N}
  < 0
\end{equation}
の利得をもたらす。
ここで$e - s + 1 < N$は区間がデータ全体より短いことを意味する。
\end{theorem}

\begin{proof}
KRIMPでは$X$の符号長は全データに対する使用頻度で決まるため、
$L_{\text{KRIMP}}(X) = -\text{usage}(X) \cdot \log_2(\text{usage}(X)/U)$
（$U$は総使用回数）である。
TCTでは区間内使用頻度が基準となり、
分母が$U$から区間内の総使用回数$U_{\text{in}} \leq U$に変わる。
区間サイズが$N$より小さいとき、
$\text{usage}(X)/U_{\text{in}} > \text{usage}(X)/U$
となるため符号長が短くなる。
\end{proof}

\subsection{計算量解析}

\begin{proposition}[計算量]\label{prop:complexity}
$|\text{Candidates}| = C$、トランザクション数$N$、
最大パターン長$\ell$のとき、
アルゴリズム\ref{alg:dic}の時間計算量は
$O(C \cdot N \cdot C \cdot \ell)$である。
\end{proposition}

各候補についてTCT全体のカバー計算（$O(N \cdot C \cdot \ell)$）を
行うため、上記の計算量となる。
実用上は$C$がMDL選択により大幅に削減されるため、
選択後のTCTサイズは小さい。


% ===================================================================
\section{実験}\label{sec:experiments}

\subsection{実験設定}

\textbf{合成データ.}
$N$個のトランザクション、$|\mathcal{I}|$個のアイテムからなる
データベースを生成する。
$P$個の密集パターン（各サイズ$|X|$）を
それぞれ長さ$30$の区間に埋め込み、
背景ノイズとして各アイテムを確率$0.1$で出現させる。

\textbf{パラメータ.}
ウィンドウサイズ$W=10$、最小サポート$\theta=3$を
デフォルト設定とする。

\textbf{評価指標.}
圧縮率$= L(\text{TCT}^*)/L(\text{SCT})$、
パターン削減率$= 1 - |\text{TCT}^*|/|\text{Candidates}|$、
実行時間を報告する。

\subsection{E1: 圧縮率評価}

表\ref{tab:e1}に様々なデータ規模での圧縮率を示す。

\begin{table}[t]
\centering
\caption{E1: 合成データにおける圧縮率}\label{tab:e1}
\begin{tabular}{ccccccc}
\toprule
$N$ & $P$ & $|X|$ & 候補数 & 選択数 & 圧縮率 & 時間(秒) \\
\midrule
100 & 2 & 2 & 11  & 2 & 0.861 & 0.003 \\
200 & 3 & 3 & 27  & 3 & 0.799 & 0.010 \\
500 & 5 & 3 & 65  & 5 & 0.879 & 0.046 \\
200 & 3 & 4 & 60  & 3 & 0.706 & 0.018 \\
\bottomrule
\end{tabular}
\end{table}

提案手法はすべての設定で圧縮を達成し、
特にパターンサイズ$|X|=4$の場合に
最も高い圧縮率（0.706、約30\%削減）を示した。
これはパターンが長いほど区間内での符号化利得が大きいことを示す。

\subsection{E2: パターン数削減率}

表\ref{tab:e2}に最大パターン長の変化に対する削減率を示す。

\begin{table}[t]
\centering
\caption{E2: MDL選択によるパターン数削減率}\label{tab:e2}
\begin{tabular}{cccc}
\toprule
最大長 & 候補数 & 選択数 & 削減率 \\
\midrule
2 & 31 & 4 & 87.1\% \\
3 & 53 & 4 & 92.5\% \\
4 & 59 & 4 & 93.2\% \\
\bottomrule
\end{tabular}
\end{table}

MDL選択は候補パターンの87--93\%を削減し、
データを簡潔に記述する少数のパターンのみを保持する。
候補数が増加しても選択数は安定しており、
MDL原理による自動的なモデル複雑度制御が機能している。

\subsection{E3: ウィンドウサイズ感度分析}

表\ref{tab:e3}にウィンドウサイズの影響を示す。

\begin{table}[t]
\centering
\caption{E3: ウィンドウサイズ感度分析}\label{tab:e3}
\begin{tabular}{cccc}
\toprule
$W$ & 候補数 & 選択数 & 圧縮率 \\
\midrule
5  & 19  & 3 & 0.801 \\
10 & 33  & 3 & 0.801 \\
15 & 58  & 4 & 0.813 \\
20 & 133 & 4 & 0.813 \\
30 & 218 & 4 & 0.812 \\
\bottomrule
\end{tabular}
\end{table}

ウィンドウサイズの増大に伴い候補数は急増するが（19$\to$218）、
MDL選択後のパターン数および圧縮率は安定している。
これはMDL原理がウィンドウサイズの変化に対して
ロバストなパターン選択を実現することを示す。


% ===================================================================
\section{考察}\label{sec:discussion}

\textbf{時間的局所性の効果.}
実験結果は、密集区間の時間的局所性が
パターンの符号化効率を向上させることを示した（定理\ref{thm:locality}）。
特にパターンサイズが大きいほど利得が大きく、
これは長いパターンほど時間的に局所化される傾向が強いためである。

\textbf{MDLによる自動パラメータ選択.}
ウィンドウサイズ$W$が変化しても選択パターン数が安定する
という結果は実用上重要である。
従来の密集区間検出では$W$と$\theta$の設定が出力サイズを直接制御するが、
提案手法ではMDL原理が追加的なフィルタリングを行うため、
パラメータ設定への感度が低い。

\textbf{計算効率.}
現在の実装は各候補追加時にカバー全体を再計算するため、
大規模データでは計算コストが問題となりうる。
SLIMのような差分更新方式の導入が今後の課題である。

\textbf{制約と今後の課題.}
本研究は合成データでの評価に限定されており、
実データでの検証が必要である。
また、密集区間の重なりが多い場合の
カバー計算の最適化も課題である。


% ===================================================================
\section{おわりに}\label{sec:conclusion}

本研究では、MDL原理と密集区間検出を統合した
パターン要約手法 Temporal Code Table を提案した。
提案手法は以下の特徴を持つ:
(1) 各パターンに時間的有効区間を付与することで
    符号化効率を向上させる;
(2) Time-Scoped Code Length により時間的局所性を活用する;
(3) 候補パターンの87--93\%を削減しつつ圧縮を達成する。

今後は実データでの大規模評価、
差分更新による計算効率の改善、
およびStreamKrimpとの統合を検討する。


% ===================================================================
\bibliographystyle{plain}

\begin{thebibliography}{10}

\bibitem{agrawal1994fast}
R.~Agrawal and R.~Srikant.
\newblock Fast algorithms for mining association rules.
\newblock In \emph{Proc.\ VLDB}, pp.\ 487--499, 1994.

\bibitem{vreeken2011krimp}
J.~Vreeken, M.~van Leeuwen, and A.~Siebes.
\newblock {KRIMP}: Mining itemsets that compress.
\newblock \emph{Data Mining and Knowledge Discovery}, 23(1):169--214, 2011.

\bibitem{smets2012slim}
K.~Smets and J.~Vreeken.
\newblock {SLIM}: Directly mining descriptive patterns.
\newblock In \emph{Proc.\ SDM}, pp.\ 236--247, 2012.

\bibitem{van2014streamkrimp}
M.~van Leeuwen and A.~Siebes.
\newblock {StreamKrimp}: Detecting Change in Data Streams.
\newblock In \emph{Proc.\ ECML PKDD}, 2008.

\bibitem{rissanen1978modeling}
J.~Rissanen.
\newblock Modeling by shortest data description.
\newblock \emph{Automatica}, 14(5):465--471, 1978.

\bibitem{grunwald2007minimum}
P.~D. Gr{\"u}nwald.
\newblock \emph{The Minimum Description Length Principle}.
\newblock MIT Press, 2007.

\bibitem{apriori_window}
著者名.
\newblock Apriori + スライディングウィンドウによる密集アイテムセット区間検出.
\newblock 技術報告, 2025.

\end{thebibliography}

\end{CJK}
\end{document}
